\chapter{Testing Strategies Derived from the Framework}\label{ch:testing}

The algebraic structure of Mol doubles as a testing oracle. Every theorem cataloged in \autoref{app:b} states a law that a compiled artifact must satisfy; each law can be discharged automatically as a property over finite input traces. This chapter turns the framework's theorems into a testing program: what the mathematics guarantees, the test harness re-checks on the artifacts, so that a dataplane build is certified against its own specification and not against a hand-written checklist. Each strategy below names the governing theorem and the class of tests it generates. The executable tools behind the computational claims are described in \autoref{app:a} and cataloged in \autoref{app:b}.

\section{Equational Law Testing (\autoref{thm:9}--\autoref{thm:12})}
The Lawvere theories of \autoref{ch:lawvere} present the ring, parser, and transport theories as sets of equations, and \ref{thm:9}--\ref{thm:12} assert that every model in Mol satisfies them. These equations are directly testable as properties: generate a finite sequence of operations, evaluate the left-hand side and the right-hand side on a state, and compare the resulting output traces and final states. The ring laws $\mathrm{push}\circ\mathrm{pop} = \mathrm{id}$ and $\mathrm{pop}\circ\mathrm{push} = \mathrm{id}$ on non-empty rings are the canonical pair; the parser laws (sequence distributes over choice, choice is associative, identity laws hold, left factoring is sound) and the transport laws (batch\_recv = recv$\otimes\cdots\otimes$recv; the interchange law) become regression suites that run on every build. Because \ref{thm:12} guarantees that adjacent push--pop pairs are removable, law testing doubles as a dead-code detector: an implementation that executes a removable pair without eliding it still satisfies the equation but violates the cost contract of \ref{thm:18}.

\section{Bisimulation and Equivalence Checking (\autoref{thm:6})}
Theorem \ref{thm:6} states that behavioral equivalence (bisimulation) is a congruence for both $\circ$ and $\otimes$, and that the quotient $\mathrm{Mol}/\!\!\approx$ is again a symmetric monoidal category. The congruence property is what makes equivalence checking compositional: to verify that $f \approx g$ it suffices to check the two machines directly, and the result extends to every context in which they are embedded. The \texttt{bisim} tool (see \autoref{app:a}) enumerates small finite machines and checks, exhaustively, that the bisimulation relation it computes is preserved by composition and tensor. On the test floor this licenses the standard substitution pattern: replace an implementation by its minimized form (\autoref{ch:coalgebra}) or its normal form (\autoref{ch:free}) and re-certify the whole dataplane by equivalence, not by re-testing behaviors. The notion goes back to \cite{park}, where bisimulation is proposed precisely as the congruence that makes such substitution sound.

\section{Model-Based Testing}
A molecule is, by definition, a stateful transformation with a transparent state space and step function; there is no hidden control flow. Every artifact therefore admits a reference model in the same vocabulary: write the simplest $(S,\mathit{step})$ that realizes the intended behavior, then test the implementation against the model over all inputs up to a bounded length and all reachable states. Because composition in Mol is composition of step functions (theorem \ref{thm:1}), models compose exactly as the artifacts do, and the model-based suite for a composite molecule is the composition of the suites of its parts. Model-based testing performs an equality check in the sense of \ref{thm:24}: the model and the artifact are two representatives of the same equivalence class, and the harness verifies the bisimulation directly.

\section{Conformance to Lawvere Theories}
Conformance testing is the operational face of \autoref{ch:lawvere}: a component that claims to implement an operation of the ring, parser, or transport theory must satisfy the theory's axioms as a conformance suite, with the axioms treated as MUST-level requirements. This is stronger than ordinary unit testing, because the axioms are closed under instantiation: any input pattern that matches the left side of an equation is a test case. \ref{thm:19} upgrades this to a decidability statement for finite signatures: the conformance problem for the whole runtime equational theory is decidable by normalization: so the conformance suite can in principle be exhaustive, not sampled. In practice the suite enumerates all ground instantiations up to a size bound and records coverage as the fraction of rewrite-normal forms exercised.

\section{Fuzzing for Memory Safety}
Fuzzing enters the framework in a precisely delimited role. Linearity (theorem \ref{thm:50}) guarantees that a well-typed hot path never allocates or deallocates, and totality (atom discipline, \autoref{ch:mol}) guarantees that every declared operation is total on its domain; fuzzing audits the boundaries of these guarantees. The fuzzer generates adversarial byte streams for parsers and rings and checks three invariants: no panic (totality), no allocation in declared hot paths (measured with an allocation counter; theorem \ref{thm:50}), and no out-of-bounds access (bitmask indexing is correct exactly when the capacity invariant of \ref{thm:48} holds). The value of fuzzing here is an audit: it cannot establish the theorems, but it can falsify the implementation's claim to instantiate them, and it feeds the differential and cost suites below.

\section{Differential Testing}
When two molecules claim the same behavior, the differential harness runs them side by side on identical input streams and compares output traces and final states. The correctness criterion is exactly the bisimulation of \ref{thm:6}, so a passing differential run is evidence of a bisimulation witness on the tested traces. The most productive instances are: (i) an optimized artifact versus its reference model (\autoref{ch:rewrite}); (ii) a minimized machine versus the original (\autoref{ch:coalgebra}); and (iii) a batch or SIMD variant versus the scalar version (\autoref{ch:operad}, \autoref{ch:linear}). Because of the congruence theorem, differential results propagate through composition: artifacts that differ only in subcomponents known equivalent need only the differing part tested.

\section{Cost-Annotated Regression (\autoref{thm:25}--\autoref{thm:27})}
The enrichment of \autoref{ch:enrich} attaches a cost vector to every molecule, and the theorems \ref{thm:25}--\ref{thm:27} fix its algebra: additive composition cost, the triangle inequality, tensor subadditivity, and the fusion bound $c (\mathrm{fused}\, g\circ f) \le c (f)+c (g)$ with strictness when the intermediate type never materializes. These are quantitative laws, and they turn into regression assertions: each CI run measures time, memory, and cache misses per pipeline and asserts (a) the cost of a composite does not exceed the sum of its parts (\autoref{thm:25}), (b) fused pipelines are no more expensive than their unfused decompositions (\autoref{thm:27}), and (c) normalization never increases cost, $c (\mathrm{NF} (M)) \le c (M)$, by \ref{thm:18}. Cost regressions are triggered by violations of the algebraic inequalities, which are stable under hardware changes in a way raw latency figures are not.

\section{Property-Based Testing of Rewrites (\autoref{thm:13}--\autoref{thm:19})}
The rewrite engine of \autoref{ch:rewrite} is a complete system for the runtime equational theory, and each of its structural theorems is testable as a property of the rewriting procedure itself. Termination (\ref{thm:13}) is tested by instrumenting every reduction with the lexicographic path measure and asserting strict descent. Local confluence (\ref{thm:14}) is tested by generating critical pairs, peaks $a \gets b \to c$ where two rules overlap, and asserting that their joins coincide; Newman's lemma (\cite{newman}) then yields global confluence, and the completion procedure is the classical one of \cite{knuthbendix}. Normal forms are tested for stability: rewriting a normal form yields itself (\ref{thm:16}, \ref{thm:17}), and applying any rule sequence to a term and then normalizing yields the same result, which is exactly the soundness statement \ref{thm:18}. Finally, the completeness statement \ref{thm:19} (two terms are equivalent iff their normal forms are equal) is the property that licenses the normal-form lookup table of \autoref{ch:free}; the property suite asserts it on all terms up to a size bound, complementing the exhaustive tool of \autoref{app:a}.

\section{State-Minimization Verification (\autoref{thm:38}--\autoref{thm:41})}
The coalgebraic minimization of \autoref{ch:coalgebra} produces, for every finite-state molecule, a minimal state space unique up to isomorphism (\ref{thm:38}), and minimization preserves behavior (\ref{thm:39}). Verification of a minimization step therefore has two clauses: the minimized artifact must be behaviorally equivalent to the original (a bisimulation check by \ref{thm:6}), and no strictly smaller state space may exist (a bound check using \ref{thm:40} on the composition's factors). Since \ref{thm:41} makes minimization compatible with equivalence: equivalent molecules have isomorphic minimizations; the test floor can canonicalize: every artifact that purports to be a pipeline stage is minimized and compared against the canonical minimized form, and any discrepancy is a specification-level bug, not a performance regression. The classical algorithms (\cite{nerode}, \cite{hopcroft}) underlie the check, and \cite{brzozowski} supplies the algebraic counterpart for the regular languages involved.

\section{Type-Level Guarantees}
The deepest tests are the ones that never run. Parametricity (\cite{wadler}) supplies free theorems: a polymorphic pipeline behaves the same on all types over which it is parametric, which the type checker enforces and the test suite uses to justify data-parallel SIMD by naturality (see \autoref{app:b}). Linearity, by \ref{thm:50}, moves the zero-allocation guarantee into the type system: a molecule that consumes its input exactly once cannot allocate in a well-typed hot path, and the ownership/borrow checker of the implementation language enforces this at compile time. The result is a layered test pyramid: type-level guarantees discharge the strongest claims without executing them; the property, differential, and cost suites discharge the algebraic laws; and fuzzing audits the boundaries. Together the layers instantiate every theorem of the catalog as an ongoing, automated obligation.
